\part{Cierre dodecafásico y dependencias transversales del estado}

\chapter[Derivación interna de la constante de estructura fina]{Derivación
interna de \texorpdfstring{\(\alpha_{\rm HMT}\)}{alpha HMT}: genealogía de
supervivencia, cierre dodecafásico, dos construcciones compatibles y
bifurcación de Witt}
\label{chap:alpha-rev7}
\HMTClaim{HMT-R-ALPHA-DODECAPHASE}
\HMTClaim{HMT-R-DODEC-EXTRACTOR}

\paragraph{Cuatro estratos de \(\alpha_{\rm HMT}\).}
\begin{enumerate}
 \item El constructor genealógico es el estado terminal enriquecido con sus
 dos publicaciones paralelas \(B_{\rm term}\) y \(U_{12}^{\rm sgn}\).
 \item La caracterización analítica es la cadena reversible
 \(U_{12}^{\rm sgn}\leftrightarrow K\to A\to\alpha_{12}\), seguida por su
 prolongación coinductiva.
 \item Al ser adimensional, su transducción común no introduce una unidad:
 fija la carta angular \(A_{\rm deg}=1000\alpha_{\rm HMT}\) utilizada por
 los canales posteriores.
 \item La comparación con CODATA se abre únicamente después de obtener
 \(\alpha_{\rm HMT}\); no selecciona \(K\), el acarreo ni la prolongación.
\end{enumerate}
La genealogía completa conserva explícitamente la cadena de entrada y alcanza
un estado terminal enriquecido. Sus dos publicaciones son paralelas:
\[
\begin{aligned}
 w_6\longrightarrow w_{30}\longrightarrow R_{36}\longrightarrow G_9
 \longrightarrow x_{\rm term}
 &\xrightarrow{\ \pi_{\rm term}\ }B_{\rm term},\\
 x_{\rm term}
 &\xrightarrow{\ R_{\rm sgn}\ }U_{12}^{\rm sgn}
 \xleftrightarrow{\ \mathcal H_{12}\ }K
 \longrightarrow A\longrightarrow\alpha_{12}.
\end{aligned}
\]
\(B_{\rm term}\) no genera \(U_{12}^{\rm sgn}\): ambos conservan
coordenadas diferentes del mismo estado. La transformación
\(\mathcal H_{12}\) es invertible, y el extractor
\((D_3K,D_4K,Q(K))\) tiene rango \(12\). En este dominio de primitivas
terminales explícitas, la cadena no recibe \(\alpha\), CODATA, una masa ni
otra coordenada metrológica como entrada.
La evaluación arquimediana forma \((P,E,\Phi,K,c_{13}=0)\), aplica el cociclo
de acarreo y deriva \(\alpha_{\rm HMT}\). La evaluación de incidencia
transporta las mismas palabras y el mismo sello a
\(B_K\subset H_\alpha^-\), que será la entrada de la rama Witt--Golay. La
construcción E21/22 y la prolongación remota
\(\mathsf T_{\alpha,\infty}\) ---denominada \(T1\) en la notación
original--- conservan dominios y mapas propios como desarrollos de la
genealogía dodecafásica principal.

\HMTDemoteHeadings
\HMTInputLive{sections/hmt/08_dodecafase_alpha}
\HMTInputLive{sections/hmt/08a_alpha_dos_vias_t1_rev8}
\HMTInputLive{sections/hmt/08b_alpha_t1_10000_rev9}
\HMTRestoreHeadings

\HMTInputLive{sections/ampliacion_20260818/03_extension_holografica_euler}

\HMTInputLive{sections/hmt/07b_realizacion_analitica_contraangulo}
\HMTInputLive{sections/hmt/16c_transportador_positivo_accion}
\HMTInputLive{sections/hmt/delta_297/16c_entrelazador_dodecafase_witt}
\begingroup
\def\HMTInternalTraceability{1}
\HMTInputLive{sections/hmt/delta_297/16d_medida_accion_positiva}
\endgroup

\HMTInputLive{sections/md/04b_banda_particula_virtual}

\chapter*{Construcción holográfica completa del electrón como dependencia del lector}
\addcontentsline{toc}{chapter}{Construcción holográfica completa del electrón como dependencia del lector}
\HMTInputLive{sections/md/06f_biografia_electron_rev6}

\chapter*{Transducción metrológica de estados genealógicos}
\addcontentsline{toc}{chapter}{Transducción metrológica de estados genealógicos}
\HMTInputLive{sections/ampliacion_20260818/07_transduccion_metrologica_estados_genealogicos}
